\section{The connected index at a collision}\label{sec:collision}

Set $g=2$. The plane Hilbert square parametrizes both pairs of distinct
points and length-two schemes supported at a single point. To extract the
connected contribution we compare its index with the symmetric square
of the degree-one index, on the space of cycles itself. All constructions below can be pulled back
over $B$; we suppress its trivial structure-sheaf factor.

Let $h:H=\operatorname{Hilb}^2(S)\to W=\operatorname{Sym}^2(S)$ be
the Hilbert--Chow morphism, $p:S^2\to W$ the finite quotient and
$i:S\to W$ the map $x\mapsto 2[x]$. Put
\begin{equation}
 K_1=\lambda_{-1}(T_S^\vee)^2,\qquad
 K_2=\lambda_{-1}(T_H^\vee)^2 .
\end{equation}
Here $\lambda_{-1}(U)=\sum_j(-1)^j[\bigwedge^j U]$. Its even and odd
summands give the signs in a permutation of two such factors. The
connected class is
\begin{equation}
 D_+=h_*K_2-\bigl(p_*(K_1\boxtimes K_1)\bigr)^{\mathfrak S_2}
                    \quad\in G_0^q(W).
 \label{eq:connected-class}
\end{equation}
The invariant summand is exact over $\mathbb C$. Over two distinct points,
$h$ is an isomorphism, and on the ordered cover the tangent bundle is the
direct sum of the tangent bundles at the two points. Exchanging them acts
on the exterior algebra with the graded permutation sign. Thus the two
terms of \eqref{eq:connected-class} restrict to the same class there.
The localization sequence for coherent sheaves makes $D_+$ the pushforward
of an integral class on $i(S)$.

Write $y=q+q^{-1}$, $e=(1-q)(1-q^{-1})=2-y$, and
$e_2=(1-q^2)(1-q^{-2})=4-y^2$. The degree-one index is $e$.
At the two ideals $(z_1^2,z_2)$ and $(z_1,z_2^2)$, in plane coordinates
$z_1,z_2$, the four tangent weights on $H$ are
$q,q^{-1},q^2,q^{-2}$. Equivariant localization divides out one
cotangent Euler factor and leaves $ee_2$ at each ideal.
Their sum is $2e^2(4-e)$, the disconnected genus-two double already
obtained from the antidiagonal Gromov--Witten evaluation
\cite[\S6]{wuebben-antidiagonal}.
For the independent pair the identity permutation has trace $e^2$;
the transposition has trace $\psi^2(e)=e_2$, where the Adams operation
$\psi^2$ replaces $q$ by $q^2$. Taking invariants averages these traces.
The connected index is therefore
\begin{align}
 F_2(q)=\chi_q(D_+)
 &=2ee_2-\tfrac12(e^2+e_2)
   =2(y-2)(y^2-3) \nonumber\\
 &=2q^{-3}-4q^{-2}+4-4q^2+2q^3 .
 \label{eq:double-character}
\end{align}
The lift to the collision diagonal is determined as well.
Equivariant homotopy invariance identifies $G_0^q(S)$ with the Laurent
character ring times $[\mathcal O_S]$. The localized Euler characteristic
of $A(q)[\mathcal O_S]$ is $A(q)/e$. Since this map is injective, the
class constructed in \eqref{eq:connected-class} is
\begin{equation}
                  D_+=i_*\big(e F_2(q)[\mathcal O_S]\big) .
\end{equation}
Each punctual ideal contributes $+1$ at $q^{-3}$ and $-2$ at
$q^{-2}$. The graded symmetric square is $(e^2+e_2)/2=4-2y$.
At $q^{-2}$ the identity and the transposition have coefficients $+1$
and $-1$, respectively, so their average vanishes. Thus the two
punctual ideals account for both leading coefficients of the connected
index. The subtraction changes the constant and the $q^{\pm1}$ terms.

For completeness one can make the connected subtraction in all degrees
at the level of characters. Introduce a curve variable $Q$ and put
$Z=1+\sum_{m\ge1}Q^m Z_m(q)$. Define integer Laurent polynomials
$F_m=\sum_j a_{m,j}q^j$ successively by
\begin{equation}
 Z=\prod_{m\ge1}\prod_{j\in\mathbb Z}(1-Q^m q^j)^{-a_{m,j}} .
 \label{eq:euler-product}
\end{equation}
After dividing by the previously determined factors, the coefficient
of $Q^m$ fixes $F_m$ integrally. Each $Z_m$ is invariant under inversion
of $q$, so uniqueness gives $F_m\in\mathbb Z[2-q-q^{-1}]$.
For $g>1$ each $Z_m$ is divisible by $e^{g-1}$, since every hook Euler
factor is divisible by $e$. The same divisibility holds for $F_m$ by
induction. Indeed the nonconstant contributions from lower degrees in
the symmetric-algebra expansion contain an Adams image of a previous
$F_d$, and $2-q^r-q^{-r}$ is divisible by $e$. This proves divisibility
over $\mathbb Q$; division by the monic polynomial $e^{g-1}$ in
$\mathbb Z[e]$ proves it integrally. For $g=1$ the hook sum gives the
ordinary partition numbers and the Euler product has $F_m=1$.

Writing $F_m=\sum_{h\ge1} n_{m,h} e^{h-1}$ and putting $q=\exp(iu)$
gives the local Gopakumar--Vafa form \cite{gv,bryanpandharipande},
\begin{equation}
 \log Z=\sum_{m,h,r\ge1} \frac{n_{m,h}}{r}
        \bigl(2\sin(ru/2)\bigr)^{2h-2} Q^{mr}.
\end{equation}
In the example \eqref{eq:double-character}, $F_2=-2e+8e^2-2e^3$.
The construction \eqref{eq:connected-class} realizes this
symmetric-algebra subtraction by an actual support calculation in
coherent $K$-theory.
